\section{结语：边疆秩序由多个国家共同承担}

辽、西夏与宋之间的长期并立说明，边界并非中心向外投射的一条线。岁币、互市、城寨、灌溉、寺院和地方官署把不同政权的生计与安全连接起来；它们既限制全面战争，也把维护秩序的成本分配给不同地区和人群。辽的多中心治理与西夏的文书、宗教和绿洲网络，不能被一概化为“边疆制度”，更不能由宋方记录替它们作结。

女真兴起和蒙古扩张改变了这些关系，却不是把旧秩序简单替换为新朝代。城市、道路、行政人员和居民会被接收、迁徙或毁坏，留下不均衡的连续性。金朝的形成正须从这个多方重组的世界理解：它接收的不只是领土，也是在辽、宋、西夏竞争中已经被深刻塑造的华北与东北社会。
